\section{Resultados}
\label{sec:resultados}
% Datos, figuras y tablas SIN interpretación. Organizado por fenómeno, no por ítem de la consigna.
% Toda figura citada por número; pies de figura autocontenidos.

\subsection{Pérdida de la asociatividad}
\label{subsec:res-asociatividad}

Con $b$ de tipo entero, las tres asociaciones de $a_N = \sum_{k=1}^{N} (1 + b^{k} - b^{k})$
devuelven exactamente $N$ para todo $N$ del rango y para las cuatro bases ensayadas. Las tres curvas
de la Figura \ref{fig:exp1-int} coinciden entre sí y con la recta $a_N = N$ en los cuatro paneles.

\begin{figure}[htbp]
  \centering
  \includegraphics[width=\textwidth]{exp1-asociaciones-int.pdf}
  \caption{Valor de $a_N = \sum_{k=1}^{N} (1 + b^{k} - b^{k})$ para $N = 1, 2, \ldots, 1000$ con las
  tres asociaciones y $b$ de tipo \texttt{int}, en escala doblemente logarítmica. Cada panel
  corresponde a una base. La línea gris continua marca el valor exacto $a_N = N$. Las tres
  asociaciones coinciden con ella en todo el rango y por eso se superponen.}
  \label{fig:exp1-int}
\end{figure}

Con $b$ de tipo flotante el resultado cambia. La Figura \ref{fig:exp1-float} muestra que $a_N$
acompaña a $N$ solo hasta un valor $k^{\ast}$ que depende de la base y después se estanca en una
meseta. Más adelante la curva desaparece del gráfico: $a_N$ vale \texttt{NaN} a partir del primer
$k$ en que $b^{k}$ desborda a infinito. La Tabla \ref{tab:exp1-umbrales} reúne los tres números que
caracterizan cada base: el $k^{\ast}$ que predice la condición $b^{k} \ge 2^{53}$ de la sección
\ref{subsec:punto-flotante}, la altura de la meseta que alcanza cada asociación y el primer $k$ con
$b^{k} = \infty$. El $k^{\ast}$ predicho coincide con la meseta medida en todas las bases y
asociaciones salvo en $b = \pm 2.0$, donde alguna asociación difiere en un término. Las dos
primeras asociaciones dan resultados idénticos en las ocho bases de la tabla. La tercera se aparta
de las otras dos solo en $b = 2.0$, donde retiene un término más (53 contra 52), y en $b = -2.0$,
donde la asimetría se invierte.

\begin{figure}[htbp]
  \centering
  \includegraphics[width=\textwidth]{exp1-asociaciones-float.pdf}
  \caption{Valor de $a_N$ para $N = 1, 2, \ldots, 1000$ con las tres asociaciones y $b$ de tipo
  \texttt{float}, en escala doblemente logarítmica. La línea gris continua marca el valor exacto
  $a_N = N$ y la línea vertical de trazo y punto, el primer $k$ en que $b^{k}$ desborda a infinito.
  Las curvas se detienen ahí porque $a_N$ vale \texttt{NaN} de ese punto en adelante. Con
  $b = 2.0$ no hay desborde dentro del rango.}
  \label{fig:exp1-float}
\end{figure}

\begin{table}[htbp]
  \centering
  \caption{Umbrales medidos por base y asociación con $b$ de tipo \texttt{float}. La columna
  $k^{\ast}$ predicho es el mayor $k$ con $|b|^{k} < 2^{53}$; las tres siguientes son la altura de
  la meseta de $a_N$, es decir la cantidad de términos que aportan 1 antes de que el sumando se
  anule. La última columna indica el primer $k$ con $|b|^{k} = \infty$.}
  \label{tab:exp1-umbrales}
  \begin{tabular}{lrrrrr}
\toprule
$b$ & $k^\ast$ predicho & $1+b^k-b^k$ & $(1+b^k)-b^k$ & $(1-b^k)+b^k$ & $k$ con $b^k=\infty$ \\
\midrule
$2.0$ & $52$ & $52$ & $52$ & $53$ & sin desborde \\
$3.0$ & $33$ & $33$ & $33$ & $33$ & $647$ \\
$5.0$ & $22$ & $22$ & $22$ & $22$ & $442$ \\
$10.0$ & $15$ & $15$ & $15$ & $15$ & $309$ \\
$-2.0$ & $52$ & $53$ & $53$ & $52$ & sin desborde \\
$-3.0$ & $33$ & $33$ & $33$ & $33$ & $647$ \\
$-5.0$ & $22$ & $22$ & $22$ & $22$ & $442$ \\
$-10.0$ & $15$ & $15$ & $15$ & $15$ & $309$ \\
\bottomrule
\end{tabular}

\end{table}

El fenómeno no aparece cuando $0 \le b \le 1$. La Tabla \ref{tab:exp1-b-hasta-uno} cubre cinco
bases de ese rango, $b \in \{0.0,\, 0.1,\, 0.5,\, 0.9,\, 1.0\}$. El desvío de un término respecto de 1
nunca supera $1.11 \times 10^{-16}$, y el error final $|a_{1000} - 1000|$ es exactamente cero en
las cinco bases.
Con bases negativas, en cambio, se reproduce lo observado con las positivas. La Figura
\ref{fig:exp1-negativas} tiene la misma forma que la Figura \ref{fig:exp1-float}. Los umbrales de
las cuatro bases negativas de la Tabla \ref{tab:exp1-umbrales} coinciden con los de $|b|$ salvo la
inversión ya señalada en $b = \pm 2.0$.

\begin{table}[htbp]
  \centering
  \caption{Comportamiento con $0 \le b \le 1$. Las tres primeras columnas dan el desvío máximo de un
  término respecto de su valor exacto 1, para cada asociación y sobre $k = 1, \ldots, 1000$; la
  última, el error absoluto de la suma completa.}
  \label{tab:exp1-b-hasta-uno}
  \begin{tabular}{lrrrr}
\toprule
$b$ & $1+b^k-b^k$ & $(1+b^k)-b^k$ & $(1-b^k)+b^k$ & $|a_{1000} - 1000|$ \\
\midrule
$0.0$ & $0$ & $0$ & $0$ & $0$ \\
$0.1$ & $1.11 \times 10^{-16}$ & $1.11 \times 10^{-16}$ & $0$ & $0$ \\
$0.5$ & $1.11 \times 10^{-16}$ & $1.11 \times 10^{-16}$ & $0$ & $0$ \\
$0.9$ & $1.11 \times 10^{-16}$ & $1.11 \times 10^{-16}$ & $0$ & $0$ \\
$1.0$ & $0$ & $0$ & $0$ & $0$ \\
\bottomrule
\end{tabular}

\end{table}

\begin{figure}[htbp]
  \centering
  \includegraphics[width=\textwidth]{exp1-bases-negativas.pdf}
  \caption{Valor de $a_N$ para $N = 1, 2, \ldots, 1000$ con bases negativas de tipo \texttt{float},
  en escala doblemente logarítmica. Cada panel corresponde a una base. La línea gris continua marca
  el valor exacto $a_N = N$ y la línea vertical de trazo y punto, el primer $k$ en que $b^{k}$
  desborda a infinito; las curvas se detienen ahí porque $a_N$ vale \texttt{NaN} de ese punto en
  adelante. Con $b = -2.0$ no hay desborde dentro del rango.}
  \label{fig:exp1-negativas}
\end{figure}

Los enteros de ancho fijo dan un resultado distinto del de los flotantes. La Tabla
\ref{tab:exp1-int64} registra, para un entero con signo de 64 bits, el primer $k$ en que $b^{k}$
sale del rango representable, el valor exacto de esa potencia y el valor que efectivamente queda
almacenado, que en tres de las cuatro bases tiene el signo cambiado. Pese a eso, $a_{1000}$ vale
1000 en las cuatro bases, igual que con los enteros de Python.

\begin{table}[htbp]
  \centering
  \caption{Desborde de los enteros de 64 bits. Para cada base se indica el primer $k$ cuya potencia
  excede el rango de un entero con signo de 64 bits, el valor exacto de $b^{k}$, el valor
  almacenado tras el desborde y el resultado de la suma completa, calculado con la simulación de
  la aritmética de 64 bits y con el entero de precisión arbitraria de Python.}
  \label{tab:exp1-int64}
  \begin{tabular}{lrrrrr}
\toprule
$b$ & $k$ del desborde & $b^k$ exacto & $b^k$ en 64 bits & $a_{1000}$ (64 bits) & $a_{1000}$ (Python) \\
\midrule
$2$ & $63$ & $9.22 \times 10^{18}$ & $-9.22 \times 10^{18}$ & $1000$ & $1000$ \\
$3$ & $40$ & $1.22 \times 10^{19}$ & $-6.29 \times 10^{18}$ & $1000$ & $1000$ \\
$5$ & $28$ & $3.73 \times 10^{19}$ & $3.59 \times 10^{17}$ & $1000$ & $1000$ \\
$10$ & $19$ & $1.00 \times 10^{19}$ & $-8.45 \times 10^{18}$ & $1000$ & $1000$ \\
\bottomrule
\end{tabular}

\end{table}

El costo de la aritmética exacta se midió aparte. La Tabla \ref{tab:exp1-costo} muestra que el
tamaño de $b^{N}$ crece linealmente con $N$, hasta 332\,193 bits (43.3 KiB) en $N = 100\,000$,
mientras que el tiempo de la corrida con enteros crece más rápido que el de flotantes: el cociente
entre ambos crece con $N$ y llega a dos órdenes de magnitud en $N = 100\,000$.

\begin{table}[htbp]
  \centering
  \caption{Costo de los enteros de precisión arbitraria con $b = 10$. Se informa el tamaño de la
  última potencia calculada, en bits y en KiB, y el tiempo de la corrida completa con \texttt{int}
  y con \texttt{float}, en milisegundos.}
  \label{tab:exp1-costo}
  \begin{tabular}{lrrrrr}
\toprule
$N$ & bits de $b^N$ & memoria de $b^N$ (KiB) & $t$ con int (ms) & $t$ con float (ms) & cociente \\
\midrule
$1\,000$ & $3\,322$ & $0.5$ & $0.23$ & $0.05$ & $4$ \\
$5\,000$ & $16\,610$ & $2.2$ & $3.50$ & $0.25$ & $14$ \\
$10\,000$ & $33\,220$ & $4.4$ & $12.41$ & $0.51$ & $25$ \\
$20\,000$ & $66\,439$ & $8.7$ & $44.46$ & $1.01$ & $44$ \\
$50\,000$ & $166\,097$ & $21.7$ & $252.51$ & $2.66$ & $95$ \\
$100\,000$ & $332\,193$ & $43.3$ & $987.15$ & $5.21$ & $189$ \\
\bottomrule
\end{tabular}


  \vspace{0.5em}
  \begin{minipage}{0.92\textwidth}
    \small \textit{Nota.} Cada tiempo es el mínimo de tres repeticiones tomadas con
    \texttt{time.perf\_counter()} y depende de la máquina y de la carga del sistema, así que tanto
    los valores absolutos como su cociente cambian entre corridas.
  \end{minipage}
\end{table}

\subsection{Dependencia del orden de suma}
\label{subsec:res-orden}

Los cuatro algoritmos de suma aplicados a $b_N = \sum_{k=1}^{N} 1/(k(k+1))$ producen resultados
distintos entre sí, y la diferencia crece con $N$. La Figura \ref{fig:exp2-error} separa las cuatro
curvas en dos grupos. Acumular de mayor a menor módulo y acumular en orden aleatorio dan errores
relativos que crecen con $N$ y llegan al orden de $10^{-14}$ en el extremo del rango grande.
Acumular de menor a mayor módulo y la suma de Kahan se mantienen en el orden de $10^{-16}$ sobre
las dos grillas, sin tendencia creciente. La cota $\sqrt{N}\,u$ de la ecuación
\eqref{eq:cota-raiz}, dibujada en la misma figura, queda por encima de las dos curvas que crecen en
todo el rango.

\begin{figure}[htbp]
  \centering
  \includegraphics[width=\textwidth]{exp2-error-relativo.pdf}
  \caption{Error relativo de $b_N$ respecto de la forma cerrada $N/(N+1)$ para los cuatro
  algoritmos de suma, en escala \texttt{symlog} en el eje vertical, que permite mostrar los valores
  exactamente nulos junto con los positivos. El panel (a) cubre $N = 10, 20, \ldots, 10\,000$ y el
  panel (b), $N = 1000, 2000, \ldots, 10^{6}$. Las líneas grises marcan la cota estadística
  $\sqrt{N}\,u$ y la unidad de redondeo $u = 2^{-53}$. Se usan marcadores en lugar de líneas porque
  las curvas de menor a mayor y de Kahan se superponen casi por completo.}
  \label{fig:exp2-error}
\end{figure}

\begin{table}[htbp]
  \centering
  \caption{Error relativo de $b_N$ respecto de $N/(N+1)$ para los cuatro algoritmos, en seis
  valores de $N$ separados por un factor 10.}
  \label{tab:exp2-errores}
  \begin{tabular}{lrrrr}
\toprule
$N$ & mayor a menor & menor a mayor & randomizada & Kahan \\
\midrule
$10$ & $0$ & $0$ & $1.22 \times 10^{-16}$ & $0$ \\
$100$ & $3.36 \times 10^{-16}$ & $0$ & $3.36 \times 10^{-16}$ & $0$ \\
$1\,000$ & $6.67 \times 10^{-16}$ & $1.11 \times 10^{-16}$ & $4.45 \times 10^{-16}$ & $0$ \\
$10\,000$ & $6.66 \times 10^{-16}$ & $0$ & $1.11 \times 10^{-15}$ & $0$ \\
$100\,000$ & $1.31 \times 10^{-14}$ & $1.11 \times 10^{-16}$ & $1.34 \times 10^{-14}$ & $1.11 \times 10^{-16}$ \\
$1\,000\,000$ & $4.76 \times 10^{-14}$ & $0$ & $1.41 \times 10^{-14}$ & $0$ \\
\bottomrule
\end{tabular}

\end{table}

La Tabla \ref{tab:exp2-errores} da los valores puntuales en seis $N$ y la Tabla
\ref{tab:exp2-resumen} resume las dos grillas completas. Sobre los mil valores de $N$ del rango
grande, la suma de menor a mayor módulo y la de Kahan alcanzan el valor de referencia sin error
medible en 599 y 862 casos respectivamente, contra 1 caso de la suma de mayor a menor y 5 de la
randomizada. Los máximos de las dos primeras se mantienen entre $1.11$ y $1.15 \times 10^{-16}$, es
decir en el nivel de $u$, y no crecen al pasar de la grilla chica a la grande. Los de las otras dos
crecen un orden de magnitud entre una grilla y la otra.

\begin{table}[htbp]
  \centering
  \caption{Resumen de los cuatro algoritmos sobre las dos grillas completas, de mil valores de $N$
  cada una. Las dos primeras columnas dan el error relativo máximo en cada grilla y la última, en
  cuántos valores de $N$ de la grilla grande el resultado coincide bit a bit con la forma cerrada.}
  \label{tab:exp2-resumen}
  \begin{tabular}{lrrr}
\toprule
algoritmo & máximo, $N \leq 10^{4}$ & máximo, $N \leq 10^{6}$ & $N$ con $E_\mathrm{rel} = 0$ \\
\midrule
mayor a menor módulo & $1.67 \times 10^{-15}$ & $5.03 \times 10^{-14}$ & $1$ de $1000$ \\
menor a mayor módulo & $1.15 \times 10^{-16}$ & $1.11 \times 10^{-16}$ & $599$ de $1000$ \\
randomizada & $4.77 \times 10^{-15}$ & $6.84 \times 10^{-14}$ & $5$ de $1000$ \\
Kahan & $1.12 \times 10^{-16}$ & $1.11 \times 10^{-16}$ & $862$ de $1000$ \\
\bottomrule
\end{tabular}

\end{table}

Repetir la suma randomizada 50 veces con $N = 10^{6}$, cada vez con una permutación distinta,
produce 48 resultados diferentes: dos pares de repeticiones coinciden. La Figura
\ref{fig:exp2-dispersion} muestra la dispersión de esos errores junto con los valores de los
algoritmos deterministas, y la Tabla \ref{tab:exp2-dispersion} la resume. El error relativo va de
$2.22 \times 10^{-16}$ a $5.64 \times 10^{-14}$, con mediana $2.10 \times 10^{-14}$ y desvío
estándar $1.47 \times 10^{-14}$. El mayor supera al menor en un factor 254. Para el
mismo $N$, la suma de mayor a menor módulo da un error de $4.76 \times 10^{-14}$, por encima de la
mediana de la randomizada y por debajo de su máximo.

\begin{figure}[htbp]
  \centering
  \includegraphics[width=\textwidth]{exp2-dispersion-randomizada.pdf}
  \caption{Dispersión de la suma randomizada con $N = 10^{6}$ sobre 50 repeticiones, cada una con
  una permutación distinta derivada de la semilla maestra. El panel izquierdo muestra el error
  relativo de cada repetición y el derecho, su histograma. Las dos líneas horizontales a trazos
  marcan los algoritmos deterministas para el mismo $N$: la de arriba, el error de la suma de mayor
  a menor módulo; la de abajo, el error nulo de las sumas de menor a mayor y de Kahan. El eje
  vertical usa escala \texttt{symlog}, que da lugar al valor nulo.}
  \label{fig:exp2-dispersion}
\end{figure}

\begin{table}[htbp]
  \centering
  \caption{Dispersión del error relativo de la suma randomizada sobre 50 repeticiones con
  $N = 10^{6}$.}
  \label{tab:exp2-dispersion}
  \begin{tabular}{lr}
\toprule
magnitud & valor \\
\midrule
repeticiones & $50$ \\
resultados $b_N$ distintos & $48$ \\
$E_\mathrm{rel}$ mínimo & $2.22 \times 10^{-16}$ \\
$E_\mathrm{rel}$ mediana & $2.10 \times 10^{-14}$ \\
$E_\mathrm{rel}$ máximo & $5.64 \times 10^{-14}$ \\
desvío estándar de $E_\mathrm{rel}$ & $1.47 \times 10^{-14}$ \\
cociente máximo sobre mínimo & $254$ \\
\bottomrule
\end{tabular}

\end{table}

\subsection{Representaciones equivalentes de una misma suma}
\label{subsec:res-representaciones}

Las dos representaciones de $b_N$, la suma de $1/(k(k+1))$ y la telescópica
$\sum (1/k - 1/(k+1))$, no producen el mismo flotante. Coinciden bit a bit en 1 de los 1001 valores
de $N$ de la grilla, únicamente en $N = 1$, y en el resto se separan por unas pocas unidades en el
último lugar, según la última columna de la Tabla \ref{tab:exp3-bn}. Ninguna de las dos se degrada
con $N$. La Figura \ref{fig:exp3-bn} muestra las dos curvas de error contenidas por la cota
$\sqrt{N}\,u$ y sin tendencia creciente. La Tabla \ref{tab:exp3-resumen} da máximos de
$1.67 \times 10^{-15}$ y $1.33 \times 10^{-15}$ sobre la grilla completa, con error exactamente
nulo en 68 y 97 valores de $N$ respectivamente.

\begin{figure}[htbp]
  \centering
  \includegraphics[width=\textwidth]{exp3-representaciones-bn.pdf}
  \caption{Error relativo de las dos representaciones de $b_N$ respecto de la forma cerrada
  $N/(N+1)$, para $N = 1, 10, 20, \ldots, 10\,000$, en escala \texttt{symlog} en el eje vertical
  para incluir los valores nulos. Las líneas grises marcan la cota estadística $\sqrt{N}\,u$ y la
  unidad de redondeo $u = 2^{-53}$.}
  \label{fig:exp3-bn}
\end{figure}

\begin{table}[htbp]
  \centering
  \caption{Error relativo de las dos representaciones de $b_N$ respecto de $N/(N+1)$ en cinco
  valores de $N$, con $b_N^{(1)} = \sum_k 1/(k(k+1))$ y $b_N^{(2)} = \sum_k (1/k - 1/(k+1))$. La
  última columna da la distancia entre los dos resultados medida en unidades del último lugar (ulp)
  del valor exacto.}
  \label{tab:exp3-bn}
  \begin{tabular}{lrrr}
\toprule
$N$ & $E_\mathrm{rel}$ de $b_N^{(1)}$ & $E_\mathrm{rel}$ de $b_N^{(2)}$ & $|b_N^{(1)} - b_N^{(2)}|$ (ulp) \\
\midrule
$1$ & $0$ & $0$ & $0$ \\
$10$ & $0$ & $1.22 \times 10^{-16}$ & $1$ \\
$100$ & $3.36 \times 10^{-16}$ & $4.49 \times 10^{-16}$ & $1$ \\
$1\,000$ & $6.67 \times 10^{-16}$ & $5.56 \times 10^{-16}$ & $1$ \\
$10\,000$ & $6.66 \times 10^{-16}$ & $3.33 \times 10^{-16}$ & $3$ \\
\bottomrule
\end{tabular}

\end{table}

\begin{table}[htbp]
  \centering
  \caption{Resumen de las dos sucesiones del tercer experimento sobre la grilla completa de 1001
  valores de $N$. Las cinco primeras filas corresponden a $b_N$, con
  $b_N^{(1)} = \sum_k 1/(k(k+1))$ y $b_N^{(2)} = \sum_k (1/k - 1/(k+1))$, medidas contra la forma
  cerrada $N/(N+1)$. La última fila corresponde a $c_N$, que no tiene forma cerrada: $D_N$ es la
  discrepancia absoluta entre sus dos representaciones y nunca se anula.}
  \label{tab:exp3-resumen}
  \begin{tabular}{lr}
\toprule
magnitud & valor \\
\midrule
$E_\mathrm{rel}$ máximo de $b_N^{(1)}$ & $1.67 \times 10^{-15}$ \\
$E_\mathrm{rel}$ máximo de $b_N^{(2)}$ & $1.33 \times 10^{-15}$ \\
$N$ con $E_\mathrm{rel} = 0$ en $b_N^{(1)}$ & $68$ de $1001$ \\
$N$ con $E_\mathrm{rel} = 0$ en $b_N^{(2)}$ & $97$ de $1001$ \\
$N$ con $b_N^{(1)} = b_N^{(2)}$ & $1$ de $1001$ \\
$N$ con $D_N = 0$ & $0$ de $1001$ \\
\bottomrule
\end{tabular}

\end{table}

La sucesión $c_N$ se comporta de otra manera. La discrepancia $D_N$ entre sus dos representaciones
no es nula en ninguno de los 1001 valores de $N$ y crece de $5.55 \times 10^{-17}$ en $N = 1$ a
$2.85 \times 10^{-11}$ en $N = 10\,000$, según la Tabla \ref{tab:exp3-cn}. En términos de cifras
significativas, los dos resultados comparten 15 en $N = 1$ y 11 en $N = 10\,000$. El panel (a) de la
Figura \ref{fig:exp3-cn} muestra esa curva junto con la suma de las discrepancias de los términos
individuales, que coincide con $D_N$ en $N = 1$, donde hay un solo término, y queda por encima en
todo el resto del rango. La misma figura incluye la cota $N^{2}u$ de la sección
\ref{subsec:cancelacion}, que acota a las dos.

\begin{table}[htbp]
  \centering
  \small
  \caption{Discrepancia entre las dos representaciones de $c_N$ en cinco valores de $N$. Se informa
  el valor calculado con la representación como cociente, la discrepancia absoluta
  $D_N = |c_N^{(1)} - c_N^{(2)}|$, su cociente con $c_N$ y la cantidad de cifras decimales
  significativas en que los dos resultados coinciden.}
  \label{tab:exp3-cn}
  \begin{tabular}{lrrrr}
\toprule
$N$ & $c_N^{(1)}$ & $D_N$ & $D_N/c_N^{(1)}$ & cifras coincidentes \\
\midrule
$1$ & $0.414213562373$ & $5.55 \times 10^{-17}$ & $1.34 \times 10^{-16}$ & $15$ \\
$10$ & $1.356033188972$ & $2.22 \times 10^{-16}$ & $1.64 \times 10^{-16}$ & $15$ \\
$100$ & $2.484679664884$ & $6.22 \times 10^{-15}$ & $2.50 \times 10^{-15}$ & $14$ \\
$1\,000$ & $3.633720211117$ & $1.79 \times 10^{-13}$ & $4.91 \times 10^{-14}$ & $13$ \\
$10\,000$ & $4.784787737052$ & $2.85 \times 10^{-11}$ & $5.95 \times 10^{-12}$ & $11$ \\
\bottomrule
\end{tabular}

\end{table}

Término a término, la discrepancia relativa entre $\sqrt{k^{2}+1} - k$ y $1/(\sqrt{k^{2}+1} + k)$
crece con $k$ conforme a la cota $2k^{2}u$ de la ecuación \eqref{eq:cota-termino}, como muestra el
panel (b) de la Figura \ref{fig:exp3-cn}. La Tabla \ref{tab:exp3-terminos} compara el $k$ que
predice esa cota con el primero que efectivamente cruza cada umbral: la discrepancia supera
$10^{-12}$ en $k = 87$, $10^{-8}$ en $k = 8193$ y $10^{-6}$ en $k = 68\,985$, entre 1.03 y 1.33
veces el valor predicho. En el último umbral la coincidencia es exacta: la cota sitúa el colapso
completo en $k = 2^{26} = 67\,108\,864$ y ese es el primer $k$ en que $\sqrt{k^{2}+1}$ devuelve
exactamente $k$. Ahí el término calculado como diferencia vale cero y el calculado como cociente
sigue siendo positivo.

\begin{figure}[htbp]
  \centering
  \includegraphics[width=\textwidth]{exp3-cancelacion-cn.pdf}
  \caption{Cancelación en $c_N$, con $c_N^{(1)} = \sum_k 1/(\sqrt{k^{2}+1} + k)$ la representación
  como cociente y $c_N^{(2)} = \sum_k (\sqrt{k^{2}+1} - k)$ la representación como diferencia. El
  panel (a) muestra la discrepancia absoluta
  $D_N = |c_N^{(1)} - c_N^{(2)}|$ entre las dos representaciones para
  $N = 1, 10, 20, \ldots, 10\,000$, junto con la suma de las discrepancias de los términos, que es
  lo que valdría $D_N$ si los errores de cada término no se compensaran entre sí, y con la cota
  $N^{2}u$. El panel (b) muestra la discrepancia relativa $\delta_k$ entre las dos formas del
  término $k$ sobre una grilla logarítmica hasta $k = 10^{8}$, con la cota $2k^{2}u$ y con el
  primer $k$ en que la resta devuelve cero. Los dos paneles usan escala logarítmica en los dos
  ejes.}
  \label{fig:exp3-cn}
\end{figure}

\begin{table}[htbp]
  \centering
  \caption{Umbrales de la discrepancia relativa $\delta_k$ entre las dos representaciones del término $k$ de
  $c_N$. Para cada umbral $\varepsilon$ se compara el $k$ que predice la cota $2k^{2}u$ con el
  primer $k$ medido que alcanza ese umbral.}
  \label{tab:exp3-terminos}
  \begin{tabular}{lrrr}
\toprule
$\varepsilon$ & $k$ que predice $2k^{2}u$ & primer $k$ con $\delta_k \ge \varepsilon$ & $\delta_k$ en ese $k$ \\
\midrule
$1 \times 10^{-12}$ & $68$ & $87$ & $1.21 \times 10^{-12}$ \\
$1 \times 10^{-10}$ & $672$ & $895$ & $1.01 \times 10^{-10}$ \\
$1 \times 10^{-8}$ & $6\,711$ & $8\,193$ & $1.12 \times 10^{-8}$ \\
$1 \times 10^{-6}$ & $67\,109$ & $68\,985$ & $1.00 \times 10^{-6}$ \\
$1$ & $67\,108\,864$ & $67\,108\,864$ & $1$ \\
\bottomrule
\end{tabular}

\end{table}

Las dos formas cerradas de $b_N$, en cambio, casi siempre coinciden. La Tabla
\ref{tab:exp3-cerradas} muestra que $1 - 1/(N+1)$ y $N/(N+1)$ difieren en 4 de los 1001 valores de
$N$ de la grilla y en 19 de los primeros 99\,999, siempre por exactamente 1 ulp. Medidos contra el
valor racional exacto, el error relativo máximo es $1.11 \times 10^{-16}$ para la primera forma y
$5.55 \times 10^{-17}$ para la segunda, es decir la mitad. Las dos saturan en 1 alrededor de
$N = 1.8 \times 10^{16}$, con un valor de $N$ de diferencia entre una y otra.

\begin{table}[htbp]
  \centering
  \caption{Comparación de las dos formas cerradas de $b_N$. Las discrepancias se cuentan sobre la
  grilla del experimento y sobre todos los $N$ menores que $100\,000$; los errores relativos se
  miden contra el valor racional exacto calculado con aritmética de fracciones. Las dos últimas
  filas dan el primer $N$ en que cada forma devuelve exactamente 1.}
  \label{tab:exp3-cerradas}
  \begin{tabular}{lr}
\toprule
magnitud & valor \\
\midrule
$N$ de la grilla donde difieren & $4$ de $1001$ \\
cuáles son esos $N$ & $170$, $510$, $1\,970$, $3\,590$ \\
$N <$ $100\,000$ donde difieren & $19$ de $99\,999$ \\
separación cuando difieren & $1$ ulp \\
$E_\mathrm{rel}$ máximo de $1 - 1/(N+1)$ & $1.11 \times 10^{-16}$ \\
$E_\mathrm{rel}$ máximo de $N/(N+1)$ & $5.55 \times 10^{-17}$ \\
primer $N$ con $1 - 1/(N+1) = 1$ & $18\,014\,398\,509\,481\,982$ \\
primer $N$ con $N/(N+1) = 1$ & $18\,014\,398\,509\,481\,983$ \\
\bottomrule
\end{tabular}

\end{table}
